\chapter{Lecture 12: Elision, Rule of Five, Const Methods and Encapsulation}
\begin{center}
\textit{February 13, 2018}
\end{center}

\section{Copy/Move Elision}
(\emph{Elision} means omission, disappearance.)
\begin{lstlisting}[style=cpp]
Vec makeVec() { return {0, 0}; }
...
Vec v = makeVec();
\end{lstlisting}
In theory: the basic ctor runs inside \code{makeVec}; then since we return by
value and the returned value is an rvalue, the move ctor runs to build \code{v}
(or the copy ctor if there is no move ctor).

What actually happens: only the basic ctor runs! The object is created directly
in the space reserved for \code{v}.
\begin{center}
\begin{tikzpicture}[>=Stealth,font=\small,o/.style={draw,rounded corners,fill=blue!5,minimum height=0.75cm,align=center}]
  \node[anchor=east,font=\scriptsize] at (-0.3,0) {in theory:};
  \node[o] (a) at (1.4,0) {\code{Vec\{0,0\}}\\ \scriptsize in \code{makeVec}};
  \node[o] (b) at (5.2,0) {\code{v}};
  \draw[->] (a) -- node[above,font=\scriptsize]{move (or copy) ctor} (b);
  \node[anchor=east,font=\scriptsize] at (-0.3,-1.3) {with elision:};
  \node[o,fill=green!10] (c) at (5.2,-1.3) {\code{v}};
  \draw[->] (1.4,-1.3) node[left=-2pt,font=\scriptsize]{basic ctor} -- (c);
\end{tikzpicture}
\end{center}

\begin{important}[Elision]
Compilers are allowed (though, in C++14, not required) to eliminate copy and
move constructions in these situations --- even if the copy/move ctor has side
effects (like printing). So never rely on the number of times a copy ctor runs.
\end{important}
To turn it off (e.g.\ to observe the calls): \code{g++ -fno-elide-constructors}.

Elision also applies to arguments:
\begin{lstlisting}[style=cpp]
void f(Vec v) { ... }
f(makeVec());    // the result is built directly in f's parameter
\end{lstlisting}

\section{Rule of Five}
\begin{important}[Rule of Five]
If you need to write any one of
\begin{itemize}
  \item the copy constructor,
  \item the copy assignment operator,
  \item the destructor,
  \item the move constructor,
  \item the move assignment operator,
\end{itemize}
then you usually need to write all five.
\end{important}
The reason: needing a custom one of these almost always means the class
\emph{owns a resource} (heap memory, a file, \dots), and every one of the five
must then handle that resource. Most classes own nothing directly and need none
of them.

\section{Short Topic: Arrays of Objects}
\begin{lstlisting}[style=cpp]
struct Vec {
    int x, y;
    Vec(int x, int y) : x{x}, y{y} {}
};

Vec vectors[3];              // stack array of 3 Vec objects
Vec *vp = new Vec[3];        // heap array of 3 Vec objects
\end{lstlisting}
Neither compiles: each element must be constructed, and there is no default
ctor. How to fix this?
\begin{enumerate}
  \item Provide a default ctor.
  \item For a stack array, use array initialization syntax:
\begin{lstlisting}[style=cpp]
Vec vectors[3] = {Vec{0, 0}, Vec{1, 3}, Vec{2, 4}};
\end{lstlisting}
  \item Use an array of pointers to objects:
\begin{lstlisting}[style=cpp]
Vec *vp[3];                    // stack array of 3 Vec pointers
Vec **vp1 = new Vec*[3];       // heap array of 3 Vec pointers
for (int i = 0; i < 3; ++i) vp1[i] = new Vec{i, i};
...
for (int i = 0; i < 3; ++i) delete vp1[i];
delete [] vp1;
\end{lstlisting}
\end{enumerate}
\begin{center}
\begin{tikzpicture}[>=Stealth,font=\small,c/.style={draw,minimum width=1.1cm,minimum height=0.65cm}]
  \node[c,fill=blue!5] (vp) at (0,0) {};
  \node[left=2pt of vp] {\code{vp1}};
  \node[font=\scriptsize] at (0,0.65) {stack};
  \node[c] (p0) at (2.5,0) {}; \node[c,right=0pt of p0] (p1) {}; \node[c,right=0pt of p1] (p2) {};
  \node[font=\scriptsize] at (3.6,0.65) {heap: array of \code{Vec*}};
  \foreach \i in {0,1,2} {
    \node[c,fill=green!8] (v\i) at (2.1+1.8*\i,-1.6) {\code{\{\i,\i\}}};
    \draw[->] (p\i.center) -- (v\i.north);
  }
  \draw[->] (vp.center) -- (p0.west);
  \node[font=\scriptsize] at (3.9,-2.3) {heap: one \code{Vec} per \code{new}};
\end{tikzpicture}
\end{center}

\section{Short Topic: Const Methods}
\begin{lstlisting}[style=cpp]
struct Student {
    int assns, mt, final;
    float grade() { return 0.4 * assns + 0.2 * mt + 0.4 * final; }
};

const Student billy{80, 50, 75};   // can't change billy's fields
billy.grade();                     // error!
\end{lstlisting}
\paragraph{Q.} Can I call \code{billy.grade()}?

No: \code{grade} did not promise not to change the fields, even though it doesn't
actually change them. To make the promise, declare the method \code{const}
(after the parameter list):
\begin{lstlisting}[style=cpp]
struct Student {
    int assns, mt, final;
    float grade() const { return 0.4 * assns + 0.2 * mt + 0.4 * final; }
};
\end{lstlisting}
Inside a const method, \code{this} has type \code{const Student*}, so the
compiler rejects any modification of a field.

\begin{important}[Rule]
\code{const} objects can only call \code{const} methods. Make every method that
doesn't modify the object \code{const}.
\end{important}

\subsection{Mutable Fields}
Suppose we want to count how many times \code{grade} is called:
\begin{lstlisting}[style=cpp]
struct Student {
    int assns, mt, final;
    int numCalls = 0;
    float grade() const {
        ++numCalls;     // error: grade is const
        return 0.4 * assns + 0.2 * mt + 0.4 * final;
    }
};
\end{lstlisting}
This won't compile. But the call counter isn't really part of the student's
logical state, so a const method should be allowed to change it. Mark the field
\code{mutable}:
\begin{lstlisting}[style=cpp]
    mutable int numCalls = 0;
\end{lstlisting}
A \code{mutable} field can be modified in a const object or by a const method.
(Distinguish \emph{physical} constness --- no bits change --- from
\emph{logical} constness --- the object's observable state doesn't change.)

\section{Invariants and Encapsulation}
\begin{lstlisting}[style=cpp]
struct Node {
    int data;
    Node *next;
    Node(int data, Node *next) : data{data}, next{next} {}
    ~Node() { delete next; }
};

Node n1{1, new Node{2, nullptr}};   // fine
Node n2{3, nullptr};                // fine
Node n3{4, &n2};                    // disaster
\end{lstlisting}
When \code{n3} goes out of scope, its dtor runs \code{delete \&n2}, deleting a
stack address: the program is likely to crash.
\begin{center}
\begin{tikzpicture}[box/.style={draw, minimum width=13mm, minimum height=7mm},>={Stealth},font=\small]
  \node[font=\scriptsize] at (-0.2,0.7) {stack};
  \node[box] (n3) at (0,0) {4 \,$\bullet$};
  \node[left=2mm of n3] {\code{n3}};
  \node[box] (n2) at (0,-1.3) {3 \,/};
  \node[left=2mm of n2] {\code{n2}};
  \draw[->,red!60!black,thick] (n3.east) to[out=0,in=0,looseness=1.6] node[right,font=\scriptsize,align=left]{\code{next} points into the stack:\\ \code{\textasciitilde Node} will \code{delete} it} (n2.east);
  \node[font=\scriptsize] at (6.2,0.7) {heap};
  \node[box] (h1) at (5.4,0) {1 \,$\bullet$};
  \node[left=2mm of h1] {\code{n1}};
  \node[box,fill=green!8] (h2) at (7.2,0) {2 \,/};
  \draw[->] (h1) -- (h2);
  \node[font=\scriptsize,align=center] at (6.3,-0.9) {fine: \code{next} owns\\ a heap Node};
\end{tikzpicture}
\end{center}

\begin{definition}
An \textbf{invariant} is a statement/assumption about a class that must be true
for the class to function correctly.
\end{definition}
\begin{itemize}
  \item Node invariant: \code{next} is either \code{nullptr} or points to a
        heap-allocated Node (that this Node owns).
  \item Stack invariant: the last item pushed is the first popped. Easy to break
        if clients can reach the underlying array and, say, sort it.
\end{itemize}
If client code has unrestricted access to the implementation, it can break
invariants, resulting in crashes and undefined behaviour.

\begin{definition}
\textbf{Encapsulation}: classes should be treated as black boxes (capsules).
Client code gets limited access through public methods (the \emph{interface}).
\end{definition}

\begin{lstlisting}[style=cpp]
struct Vec {
    Vec(int x, int y);        // public: struct members are public by default
  private:
    int x, y;
  public:
    Vec operator+(const Vec &v) const { return {x + v.x, y + v.y}; }
};

int main() {
    Vec v{1, 2};              // OK: Vec(int, int) is public
    Vec v1 = v + v;           // OK: operator+ is public; free Big 5 are public
    cout << v.x << endl;      // won't compile: x is private
}
\end{lstlisting}

\begin{important}[Advice]
Fields should be private. Only selected methods should be public (the
interface).
\end{important}

\subsection{The \code{class} Keyword}
\begin{lstlisting}[style=cpp]
class Vec {
    int x, y;                 // private by default
  public:
    Vec(int x, int y);
    Vec operator+(const Vec &v) const;
};
\end{lstlisting}
The \emph{only} difference between \code{struct} and \code{class} is the default
visibility: public for \code{struct}, private for \code{class}.

\section{Encapsulating Node}
To guarantee the Node invariant, create a wrapper class \code{List} that has
exclusive access to the Node objects:
\begin{lstlisting}[style=cpp]
// list.h
class List {
    struct Node;                 // private nested class: declared only
    Node *theList = nullptr;
  public:
    void addToFront(int n);
    int ith(int i) const;
    ~List();
};

// list.cc
#include "list.h"
struct List::Node {              // full definition hidden in the .cc file
    int data;
    Node *next;
    Node(int data, Node *next) : data{data}, next{next} {}
    ~Node() { delete next; }
};

List::~List() { delete theList; }

void List::addToFront(int n) {
    theList = new Node{n, theList};
}

// requires: 0 <= i < length of the list
int List::ith(int i) const {
    Node *curr = theList;
    for (int j = 0; j < i && curr; ++j, curr = curr->next);
    return curr->data;
}
\end{lstlisting}
No one outside \code{List} knows the Node structure (it is defined only in the
\code{.cc} file), so only \code{List} can create Nodes, and it always puts them
on the heap. We now have a guarantee that the invariant is maintained. (We should
also write or disable the copy/move operations of \code{List}: Rule of Five.)

\begin{center}
\begin{tikzpicture}[box/.style={draw, minimum width=11mm, minimum height=7mm},>={Stealth},font=\small]
  \node[draw,fill=blue!5,minimum width=1.8cm,minimum height=0.8cm] (l) at (0,0) {\code{theList}};
  \node[font=\scriptsize,anchor=south] at (l.north) {\code{List l}};
  \node[box] (a) at (3,0) {3 \,$\bullet$};
  \node[box] (b) at (4.6,0) {2 \,$\bullet$};
  \node[box] (c) at (6.2,0) {1 \,/};
  \draw[->] (l) -- (a); \draw[->] (a) -- (b); \draw[->] (b) -- (c);
  \node[draw,dashed,rounded corners,fit=(a)(b)(c),inner sep=8pt,label={[font=\scriptsize]below:\code{List::Node}: invisible to clients}] {};
  \node[font=\scriptsize,align=center] (cl) at (-3.8,0) {client code};
  \draw[->] (cl) -- node[above,font=\scriptsize]{public methods} (l);
\end{tikzpicture}
\end{center}

\paragraph{The cost.} The client can no longer walk the list. Printing it
requires \code{ith(0)}, \code{ith(1)}, \dots, each of which walks from the
front, so traversing the list costs
$1 + 2 + \cdots + n \in \Theta(n^2)$ instead of $\Theta(n)$. Next lecture:
the Iterator pattern fixes this.

\section*{Summary}
\begin{itemize}
  \item Compilers may elide copies/moves; don't rely on their side effects.
  \item Rule of Five: write all five or none.
  \item Arrays of objects need a default ctor, an initializer, or pointers.
  \item \code{const} methods can be called on const objects; \code{mutable}
        fields are exempt from constness.
  \item Invariants must be protected by encapsulation: private fields, public
        interface; \code{class} defaults to private.
  \item Hiding Node inside List secures the invariant but makes traversal
        $\Theta(n^2)$.
\end{itemize}
